\epigraph{It is wrong to think that the task of physics is to find out how nature is. Physics concerns what we can say about Nature.}{Niels Bohr, \cite{petersenPhilosophyNielsBohr1963}}

The aim of the quantum reconstruction program is to derive quantum mechanics from accessible postulates instead of taking Hilbert space as a primitive structure \cite{hardyQuantumTheoryFive2001}. Reconstructions uncover why quantum theory assumes its characteristic form, addressing Wheeler's call to explain "Why the quantum?'', whereas interpretational approaches take the mathematical structure as given. Von Neumann's Hilbert space axiomatization supplies a powerful formalism but leaves most foundational questions unanswered \cite{vonneumannMathematicalFoundationsQuantum2018}. This drove many scientists toward operationalism, which seeks the minimal structure encapsulating laboratory operations. This approach was started by Peres \cite{peresQuantumTheoryConcepts2002} but we can track its roots from Bohr's epistemological philosophy that quantum formalism expresses communicable outcomes of macroscopic devices. Constructing the theory on these procedures selects bounded observables as primitive operational objects. At the same time, the algebraic formulation arose from Segal's observation that not every Hilbert space operator corresponds to a measurable quantity \cite{segalIrreducibleRepresentationsOperator1947}. In fact, Segal identified the $C^*$-algebra studied by Gelfand and Naimark \cite{gelfandImbeddingNormedRings1994} as the appropriate setting. Gelfand-Naimark-Segal construction associates a Hilbert space representation to each state without assuming it beforehand, deriving the familiar formalism from the operational content of the algebra's relations.

An alternative approach used to understand quantum mechanics relies on finding analogies between the quantum formalism and the classical one. Kibble discovered that a projective Hilbert space supplies a symplectic structure whose Hamiltonian flow generates Schrödinger evolution \cite{kibbleGeometrizationQuantumMechanics1979}. Then, Ashtekar and Schilling showed that there is a full Kähler structure that embraces a symplectic form, the Fubini-Study metric, and a compatible complex structure \cite{ashtekarGeometricalFormulationQuantum1997}. To conclude, Brody and Hughston formalised that expectation values and quantum uncertainties admit geometric expressions on a Kähler manifold \cite{brodyGeometricQuantumMechanics2001}. This formulation follows from the inner-product geometry, providing a unified description where dynamics, symmetries, and measurements arise from the manifold's geometry without assuming separately the Schrödinger equation. 

Motivated by this geometric perspective, the central problem of this thesis is the investigation of the structural correspondence between classical and quantum mechanics. In fact, both frameworks share multiple Hamiltonian structures. Classical mechanics evolves on a symplectic manifold $(\mathcal{M},\omega)$ via Hamiltonian vector fields generated by smooth functions. Analogously, the projective Hilbert space of quantum pure states carries a symplectic form whose associated Hamiltonian flow generates the Schr\"odinger evolution \cite{kibbleGeometrizationQuantumMechanics1979}. Another correspondence appears in the Poisson bracket to commutator analogy first recognised by Dirac when he proposed that the quantum commutator corresponds to $i \hbar$ times the classical Poisson bracket.  Moreover, Noether's theorem in both theories associates continuous symmetries with conserved quantities, with the quantum geometric version using a momentum map on the K\"ahler state space \cite{ashtekarGeometricalFormulationQuantum1997}. These parallels raise whether these are distinct theories or two realisations of a common framework. This thesis reconstructs both from a shared algebraic starting point examining where their geometries diverge.

Understanding shared algebraic and geometric structures is necessary for several reasons. Within this framework, the problem of quantization reduces to deforming a commutative classical observable algebra into a noncommutative one, with the Poisson bracket arising as the first-order term in the commutator expansion. This approach formalizes also the classical limit, that requires a rigorous framework such as Landsman's continuous fields of $C^*$-algebras in \cite{landsmanMathematicalTopicsClassical1998}. Moreover, formulating both theories within a shared algebraic language help us understand better the boundary between a quantum system and a classical measuring apparatus. 

The literature contains many classical quantum correspondences. Dirac's canonical quantization suffers from ordering ambiguity and the Groenewold van Hove obstruction \cite{groenewoldPrinciplesElementaryQuantum1946}, while Weyl's systematic map from phase space functions to operators does not preserve the Poisson brackets \cite{weylTheoryGroupsQuantum2009}. For these reasons, we present Segal's algebraic formulation that uses $C^*$-algebras treating classical and quantum mechanics on the same footing \cite{segalIrreducibleRepresentationsOperator1947}. In addition, Strocchi showed that the GNS construction yields a symplectic structure on the state space \cite{strocchiComplexCoordinatesQuantum1966} while Landsman provided the most comprehensive algebraic treatment of the classical limit \cite{landsmanMathematicalTopicsClassical1998}. Separately, the geometric formulation, initiated by Kibble, showed that a projective Hilbert space carries a symplectic form whose Hamiltonian flow is Schr\"odinger evolution \cite{kibbleGeometrizationQuantumMechanics1979}. While Heslot extended this by proving that the geometry of projective space determines the quantum formalism \cite{heslotQuantumMechanicsClassical1985}. Ashtekar and Schilling completed the picture with a canonical Kähler structure combining symplectic form, Fubini Study metric, and complex structure, with the metric's geodesic distance matching Born rule transition probabilities \cite{ashtekarGeometricalFormulationQuantum1997}. Moreover, Cirelli, Manià and Pizzocchero showed that expectation values, uncertainties, and the commutator structure have direct algebraic expressions on the state manifold \cite{cirelliQuantumMechanicsInfinitedimensional1990}. These approaches have achieved significant successes including the Poisson bracket commutator link, the unified C*-algebraic framework, and the Kähler structure. However the canonical quantization is not unique. The algebraic formulation does not explain why quantum mechanics requires noncommutativity or how the classical limit operates operationally. 

\vspace{0.5cm}
\noindent A synopsis of the thesis is the following.

Chapter one builds the classical picture from the Gelfand-Naimark theorem. This identifies a unital commutative $C^*$-algebra with continuous functions on its spectrum, so phase space is recovered as a topological space. Since that recovery gives only topology, the chapter introduces a dense smooth subalgebra with a Poisson bracket. This extra structure is not intrinsic to the $C^*$-algebra but it is needed to give the spectrum a differentiable symplectic structure. We define states as normalized positive functionals while their probabilistic interpretation comes from the Riesz-Markov theorem. In addition, we state that time evolution can be algebraically represented by a one-parameter automorphism group. A realization of this formalism is presented as the classical spin system derived from the group $\mathfrak{so}(3)$. Its Poisson bracket, via Gelfand's theorem, yields the sphere $\mathbb{S}^2$ as the phase space, where the symplectic form, the $SO(3)$ moment map, and angular momentum conservation along the symmetry axis are built.

Chapter two turns to the algebraic formulation of quantum mechanics using a non-commutative observable algebra. The GNS construction shows that every state on a $C^*$-algebra yields a Hilbert space and a representation by bounded operators, with the state recovered as an expectation on a cyclic vector. In the non-commutative case, GNS serves the same purpose as Gelfand's theorem in the commutative one. It moves from an abstract algebra to a Hilbert space representation without assuming it beforehand. It also connects pure states to irreducible representations. Because only bounded operators can lie in a $C^*$-algebra, we introduce the Weyl algebra to handle the unboundedness of position and momentum. Consequently, we show that representing the Weyl $C^*$-algebra leads to Schr\"odinger wave functions. The chapter ends with a systematic comparison of the two reconstructions. In the non-commutative setting, spectrum points are replaced by Hilbert space representations, and the Poisson bracket, restricted to the smooth subalgebra, is replaced by a Lie bracket on the whole algebra. The two-level system based on $\mathfrak{su}(2)$ provides an example. GNS applied to the $SO(3)$-invariant state yields the Pauli algebra as an irreducible representation. Comparing the Bloch ball with the simplex of classical states on the sphere shows how the same Lie algebra generates different state geometries depending on whether it uses a Poisson bracket or a commutator.

Chapter three gives a geometric reformulation of quantum mechanics. It starts with the pure state space, identified with projective Hilbert space, without assuming the Schrödinger equation or symmetries as extra ingredients. We show explicitly that this projective space carries a Kähler structure combining a symplectic form, a Riemannian metric, and a complex structure. As in classical mechanics, the symplectic form defines a Poisson bracket that drives time evolution. In fact, given the energy selfadjoint operator $H$, its expectation value is a smooth function on the state space, while the corresponding Hamiltonian vector field generates a flow whose integral curves satisfy the Schrödinger equation. The Riemannian metric accounts for statistical and probabilistic features defining a distance between pure states whose geodesic separation matches the Born's rule transition probability. Continuous symmetries appear as flows that preserve the Kähler structure, and the geometry provides a moment map pairing each symmetry generator with a conserved quantity, a geometric Noether theorem. Dynamics and symmetries all follow from the state manifold's geometry, providing a geometric formulation of quantum mechanics. In an example, we extend the previous two-level quantum system to a Kähler manifold showing the dynamical analogies with the first classical example.
\label{cap:introduzione}
